\paragraph{The group-recovery step.}
\label{sec:recovery-theory}
The recovery step scores a \emph{fit} $F$, a set of lines $\theta_k$ with scales $s_k$ and
proportions $\pi_k$ and a noise weight $\pi_0>0$ summing to one, by
$\ell_n(F)=\sum_{i\le n}\log f_F(u_i)$, where
$f_F(u)=\sum_{k}(\pi_k/s_k)\,\varphi\{(y-x^{\top}\theta_k)/s_k\}+\pi_0/R$. For given
parameters, $\ell_n(F)/n$ converges almost surely to $L(F)=\E_P\log f_F(u)$; the proposition
values, in $L$, the two operations by which a candidate differs from the original fit.

\begin{proposition}[informal statement; the precise version is in Supplement~S2.3]
\label{prop:recover-value}
(i) Let $G$ be an event of probability $\pi_\ast$ on which $y=x^{\top}\theta_\ast+\sigma_\ast e$
with $e\sim\mathcal N(0,1)$ independent of $x$, let $F$ have noise weight $\pi_0>\pi_\ast$, and
let $F'$ be $F$ with the line $(\theta_\ast,\sigma_\ast,\pi_\ast)$ added and the noise weight
lowered to $\pi_0-\pi_\ast$. Then
\begin{equation}
L(F')-L(F)\;\ge\;\pi_\ast\Bigl\{\log\frac{\pi_\ast R}{\pi_0}-\log\sigma_\ast
-\tfrac12\log(2\pi e)-\xi\Bigr\}-\nu\log\frac{\pi_0}{\pi_0-\pi_\ast},
\label{eq:recover-gain}
\end{equation}
where $\xi\ge0$ is small when $G$ lies a few scales from every line of $F$, and
$\nu\le1-\pi_\ast$ is the noise mass that $F$ attributes to the units outside $G$. (ii) If $F$
contains two lines through one group of probability $\pi_S$ and scale $\sigma_S$, at intercepts
$\delta$ apart, and $F''$ replaces them by the line through the group, then
$|L(F'')-L(F)|\le(\varphi(1)a^{2}+\tfrac14a^{4})\pi_S$ with $a=\delta/(2\sigma_S)$.
\end{proposition}

The leading term of \eqref{eq:recover-gain} is the value of a group: per unit, the log of
the ratio of its peak density $\pi_\ast/\sigma_\ast$ to the noise density $\pi_0/R$, less the
entropy $\tfrac12\log(2\pi e)\approx1.42$ of a unit normal error. With $R/\sigma_\ast=40$,
$\pi_\ast=0.1$, $\pi_0=0.2$ and $\xi\approx0$, the value is $1.58$ nats per unit of the group,
that is $0.158n$ in all. With these values the loss term is $\nu\log2$: for a fit whose noise
mass off $G$ is about $0.1$, the weight $\pi_0=0.2$ less the group's $0.1$, it is about $0.069n$
and the gain at least $0.09n$ against a margin $\tfrac{d+1}{2}\log n$ of $8.6$ for $d=2$,
$n=300$. The bound $\nu\le1-\pi_\ast$ alone gives a loss of up to $0.62n$, which exceeds the
value, so the inequality is informative only for such a fit. Part (ii), an idealisation with two parallel lines at intercepts $\delta$ apart, says that
removing one of two lines through one group is nearly free, so a candidate that arises by the two operations is accepted
for all large $n$ when the right side of \eqref{eq:recover-gain} exceeds $\varphi(1)a^{2}\pi_S$
and its scales are admissible. The result is for fixed parameters, which the step reads off
the data, and it does not cover the reweighting refit of the candidate. The peak test, analysed
in Supplement~S2.3, fails a band of outliers that fills its window and passes one of half-width
below about $3.4$ scales.
